\documentclass[12pt]{article}
\usepackage[shortlabels]{enumitem}
\usepackage{float}
\usepackage{xcolor}
\usepackage[a4paper, top=0.6cm, bottom=0.6cm, left=1.5cm, right=1.5cm]{geometry}
\usepackage{parskip}
\usepackage{setspace}
\usepackage{lmodern}
\usepackage[T1]{fontenc}
\usepackage[brazil]{babel}
\usepackage{hyperref}
\usepackage{graphicx}
\usepackage{amsmath}
\usepackage{amssymb}
\usepackage{amsfonts}
\usepackage{dsfont}

% --- Ajuste fino do título ---
\makeatletter
\renewcommand{\maketitle}{
  \begin{center}
    {\bfseries \@title \par}
  \end{center}
  \vspace{-0.1cm} % diminui o espaço entre título e texto
}
\makeatother

\title{Gabarito - Prova 1 - Estatística para Administração - 2025.2}
\date{}

\begin{document}
\maketitle


1) 

\begin{enumerate}[a)]
    \item 
    
    Denotemos por $\bar{z}$ e $\sigma_z$ a média e o desvio padrão do conjunto  $z_1, z_2, \dots, z_n$,  respectivamente. A partir disso:

    \begin{align*}
        \bar{z} &= \dfrac{\displaystyle\sum_{i=1}^n  z_i}{n}
         = \dfrac{\displaystyle\sum_{i=1}^n \frac{x_i - \bar{x}}{\sigma_x}}{n}
        = \displaystyle\sum_{i=1}^n \frac{x_i - \bar{x}}{n \sigma_x}
        = \dfrac{1}{\sigma_x} \dfrac{\displaystyle\sum_{i=1}^n (x_i - \bar{x})}{n}
        = \dfrac{1}{\sigma_x} \left(\dfrac{\displaystyle\sum_{i=1}^n x_i}{n} - \dfrac{\displaystyle\sum_{i=1}^n \bar{x}}{n}\right)\\
        &= \dfrac{1}{\sigma_x} \left(\bar{x}- \dfrac{n \bar{x}}{n}\right) 
        =\dfrac{1}{\sigma_x} \left(\bar{x} - \bar{x}\right) = 0
    \end{align*}

    \begin{align*}
        \sigma_z^2 = \displaystyle\sum_{i=1}^n \dfrac{(z_i - \bar{z})^2}{n} = \displaystyle\sum_{i=1}^n \dfrac{z_i^2}{n}
        = \displaystyle\sum_{i=1}^n \dfrac{(\frac{x_i - \bar{x}}{\sigma_x})^2}{n}
        = \dfrac{1}{\sigma_x^2}\displaystyle\sum_{i=1}^n \dfrac{(x_i - \bar{x})^2}{n}
        = \dfrac{1}{\sigma_x^2} \cdot \sigma_x^2 = 1
    \end{align*}

    \item

Denotemos por $z_m$ o escore padronizado do número de tarefas concluídas por Mariana com respeito ao conjunto de dados formado pelos estagiários da empresa e por 
$z_j$ o escore padronizado do número de tarefas concluídas por Jéssica com respeito ao setor de Marketing. Dessa forma, temos que:


$$z_m = \dfrac{3 - 2,8}{\sqrt{0,25}} = \dfrac{0,2}{0,5} = 0,4$$

$$z_j = \dfrac{8 - 7}{2,5} = \dfrac{1}{2,5} = 0,4$$

A partir disso, podemos ver que quando comparamos o escore padronizado do número de tarefas realizadas por Mariana com base no conjunto de dados formados pelos estagiários da empresa, verificamos
que o desempenho dela foi 0,4 desvios padrão acima da média. Além disso, quando comparamos o escore padronizado do número de tarefas concluídas por
Jéssica com o dos demais funcionários do setor de Marketing, também obtemos que seu desempenho foi 0,4 desvios padrão acima da média. Dessa forma, 
se olhamos para o conjunto dos estagiários, Mariana apresenta o mesmo comportamento de Jéssica, que foi o exemplo dado pelo time de Gestão de Pessoas. Com isso,
podemos concluir que a hipótese de Mariana está correta, ou seja, a equipe de Gestão de Pessoas está analisando seu desempenho com base em funcionários 
mais experientes e não com pessoas que ocupam o mesmo cargo que Mariana. 
\end{enumerate}

2) 


\begin{enumerate}[a)]
    \item 
    Definamos o evento $A:$ Cliente ter sido captado pela Quarto Andar-Edu 

    $$\mathds{P}(A) = \dfrac{6800}{18000} = \dfrac{68}{180}=\dfrac{17}{45}$$
    \item 
    Definamos o evento $B:$ O imóvel anunciado pelo cliente é uma Casa
    $$\mathds{P}(A|B) = \dfrac{1800}{9000} = \dfrac{18}{90}$$
    \item 
    Definamos o evento $C:$ O imóvel anunciado pelo cliente é um apartamento. 
    $$\mathds{P}(A|C) = \dfrac{5000}{9000} = \dfrac{50}{90}$$
    \item 
    É possível perceber que a probabilidade de um cliente ter vindo da plataforma Quarto Andar-Edu dado que o cliente possui um apartamento é maior do que 
    a de ele vir da plataforma dado que ele possui uma casa. Dessa forma, conclui-se que a plataforma parece ter sido mais efetiva em donos
    de apartamentos. Esse resultado pode indicar que o público atingido por meio da Quarto-Andar Edu é formado majoritariamente por donos de apartamentos,
    podendo gerar iniciativas dentro da empresa para segmentar a plataforma.  
    \item 
    Usamos a interpretação frequentista. A interpretação frequentista assume que a probabilidade de um evento ocorrer é igual
    a sua frequência de ocorrência quando o número de repetições (ou amostras) é muito grande. As principais críticas feitas à abordagem frequentista estão ligadas 
    à determinação do número de repetições necessárias para que a frequência seja igual a probabilidade, ademais, também notamos críticas
    com respeito a eventos que não podem ser repetidos, implicando em ser uma abordagem que só funciona em alguns casos, mas em outros não, o que 
    a tornaria de certa forma fraca. 
    \item Definamos o espaço amostral do experimento:
    $$\Omega = \{Casa, Apartamento\}^3$$
    
    Como a questão fala para usarmos a abordagem clássica, consideraremos que os clientes possuem igual probabilidade de anunciarem uma Casa
    ou um Apartamento. 

    Definamos o evento $D$: Pelo menos dois clientes anunciarem um apartamento. Escrevendo $D$ por meio da notação de conjunto, temos que:

    $$D=\{(A, A, C), (A, C, A), (C, A, A), (A, A, A)\}$$

    Ainda usando a interpretação clássica da Probabilidade, temos que:

    $$\mathds{P}(D) = \dfrac{|D|}{|\Omega|} = \dfrac{4}{8} = \dfrac{1}{2}$$
\end{enumerate}


\vspace{5px}

3) 

Definamos os seguintes eventos:

\begin{itemize}
    \item $A$: A pessoa clica no anúncio
    \item $B$: A pessoa anuncia no site da empresa. 
\end{itemize}

\begin{align*}
    \mathds{P}(A|B) &= \dfrac{\mathds{P}(B|A)\mathds{P}(A)}{\mathds{P}(B)} = \dfrac{0,95 \cdot 0,01}{\mathds{P}(B|A)\mathds{P}(A) + \mathds{P}(B|A^c)\mathds{P}(A^c) }
    = \dfrac{0,95 \cdot 0,01}{0,95 \cdot 0,01 + 0,1 \cdot 0,99}\\
    &= \dfrac{0,95}{10,85} = 0,0875
\end{align*}

Verificamos que apesar da grande proporção de clientes que anunciam na plataforma vindos do anúncio, o fato das pessoas clicarem no anúncio apenas 1\% das vezes faz com que a probabilidade de que 
um cliente da plataforma tenha vindo por meio desse anúncio seja baixa. Nesse sentido, talvez a equipe de Marketing devesse investir em melhorar o anúncio 
para aumentar o percentual de clientes que clicam, podendo talvez melhorar o design dos anúncios ou a abordagem realizada. 

\vspace{5px}

4) Sejam $A$ e $B$ dois eventos em um mesmo espaço amostral $\Omega$. Mostre matematicamente os resultados das questões $a$ e $b$ e com palavras os da $c$, $d$ e $e$:
\begin{enumerate}[a)]
    \item 
    $\mathds{P}(A|B) = \dfrac{\mathds{P}(A \cap B)}{\mathds{P}(B)}$ implica que $\mathds{P}(B|A) = \dfrac{\mathds{P}(B \cap A)}{\mathds{P}(A)}$
    além disso, garante também que  $\mathds{P}(B \cap A) = \mathds{P}(B|A) \mathds{P}(A)$

    Usando essas informações e o fato de que $A \cap B = B \cap A$, segue que:

    $$\mathds{P}(A|B) = \dfrac{\mathds{P}(A \cap B)}{\mathds{P}(B)} = \dfrac{\mathds{P}(B \cap A)}{\mathds{P}(B)} = 
    \dfrac{\mathds{P}(B|A) \mathds{P}(A)}{\mathds{P}(B)}
    $$
    
    \item 
    
    Podemos escrever $B$ como a união de $B \cap A$ e $B\cap A^c$, ou seja, $B = (B \cap A) \cup (B \cap A^c)$

Note que na relação acima funcionaria mesmo se $A \not\subset B$

A partir disso, temos:

$$\mathds{P}(B) = \mathds{P}((B \cap A) \cup (B \cap A^c)) =  \mathds{P}(B \cap A) +  \mathds{P}(B \cap A^c)$$

Obs: Note que $B \cap A$ e $B \cap A^c$ são mutuamente exclusivos.

Ademais, se $A \subset B$, então $A \cap B = A$

Como $\mathds{P}\left(B \cap A^c\right) \geq 0$, segue que $\mathds{P}(B) \geq  \mathds{P}(A \cap B)$ e, portanto, $\mathds{P}(B) \geq  \mathds{P}(A)$

    \item 

    Aqui bastava explicar que estamos calculando a probabilidade da união de eventos. Se eles não forem mutuamente exclusivos, 
    somar as probabilidades de $A$ e $B$ incorre em somar duas vezes a probabilidade de $A \cap B$, logo, precisamos descontar essa segunda vez que aparece. 
    Quando os eventos são mutuamente exclusivos, a relação ainda é válida, mas sabemos que a interseção é o vazio, que tem probabilidade zero. 
    Para ganhar o total de pontos da questão também seria necessário desenhar o Diagrama de Venn.

    \item 
    
    O fato de $A$ e $B$ serem mutuamente exclusivos implica que não podem ocorrer ao mesmo tempo. Logo, se um deles ocorre, o outro automaticamente não pode ocorrer. 

    \item 
    A resposta para essa questão consistia em  fazer o Diagrama de Venn contendo $A$ e as partições seguido de uma breve explicação do conceito de partição de um espaço amostral, citando
    que uma partição de um espaço amostral nada mais é que uma sequência de eventos mutuamente exclusivos onde a união de todos eles constitui o próprio espaço amostral. 
    Por serem mutuamente exclusivos, cada um dos conjuntos $B_i$ terá ou não uma parte de $A$, sendo essa parte não partilhada com as demais $B_j, j\neq i$, o que implica
    que os eventos $A \cap B_i, i \in \{1, \dots, n\}$ também são mutuamente exclusivos. 

    Dessa forma, temos que $A = \displaystyle\bigcup_{i=1}^n A \cap B_i$ e, portanto,:

    $$\mathds{P}(A) = \mathds{P}\left( \bigcup_{i=1}^n A \cap B_i \right) =\sum_{i=1}^n \mathds{P}(A \cap B_i)$$

    Usando o resultado da questão $a$, temos que:

    $$\mathds{P}(A) = \sum_{i=1}^n \mathds{P}(A |B_i) \mathds{P}(B_i)$$

    
\end{enumerate}

\end{document}